\section{The absorption criterion}\label{sec:absorption}

The passage from a lifted correction to an absorption map is elementary
and independent of the coordinate construction. We state it first to
specify exactly what that construction must supply.
Triangular changes of variables also appear in the quotient-section
construction of Aharoni and Lindenstrauss \cite[pp.~281--282]{AL78} and in the
free-space splitting of Godefroy and Kalton \cite[Theorem~2.12]{GK03}.
The following argument lifts a correction term and supplies its own
inverse formula.

\begin{lemma}[Absorption from a lifted correction]\label{lem:absorption}
Let \(E,U,Z\) be real Banach spaces, let
\(D=E\oplus_\infty U\), and let \(P(t,u)=u\). Suppose
\(Q:Z\to U\) is bounded linear,
\(g:D\to U\) is a bi-Lipschitz bijection, and \(K:D\to Z\) is
Lipschitz. Assume the exact identity
\begin{equation}
                 QK(t,u)=g(t,u)-u
                 \qquad ((t,u)\in D).                    \label{eq:absorb-identity}
\end{equation}
Then
\[
 F:Z\oplus_\infty E\longrightarrow Z,\qquad
 F(b,t)=b+K(t,Qb)
\]
is a bi-Lipschitz bijection. If
\[
 q=\|Q\|,\qquad L_K=\Lip(K),\qquad
 \|g(x)-g(x')\|_U\geq c_g\|x-x'\|_D
\]
for some arbitrary lower constant \(c_g>0\), then
\begin{equation}
 \Lip(F)\leq 1+L_K\max\{1,q\},\qquad
 \Lip(F^{-1})\leq
 \max\left\{1+\frac{L_Kq}{c_g},\frac{q}{c_g}\right\}.     \label{eq:absorb-bounds}
\end{equation}
No surjectivity or right-inverse assumption is imposed on \(Q\).
\end{lemma}

\begin{proof}
For \(d\in Z\), define
\begin{equation}
 (t_d,u_d)=g^{-1}(Qd),\qquad
 G(d)=\bigl(d-K(t_d,u_d),t_d\bigr).                        \label{eq:absorb-inverse}
\end{equation}
This is defined for every \(d\) because \(g\) is onto \(U\).
Identity \eqref{eq:absorb-identity} gives
\[
 QF(b,t)=Qb+QK(t,Qb)=g(t,Qb).
\]
Thus \(g^{-1}\) reads \((t,Qb)\) from \(QF(b,t)\), and \eqref{eq:absorb-inverse} gives
\(G(F(b,t))=(b,t)\). Conversely, write \((t,u)=g^{-1}(Qd)\)
and \(b=d-K(t,u)\). Then
\[
 Qb=Qd-QK(t,u)=g(t,u)-\bigl(g(t,u)-u\bigr)=u,
\]
so \(F(G(d))=d-K(t,u)+K(t,Qb)=d\). This proves bijectivity with
inverse \(G\).

The map \((b,t)\mapsto(t,Qb)\) has Lipschitz constant at most
\(\max\{1,q\}\), which proves the first estimate in \eqref{eq:absorb-bounds}.
For \(d,d'\in Z\), the lower bound for \(g\) gives
\[
 \|(t_d,u_d)-(t_{d'},u_{d'})\|_D
       \leq \frac{q}{c_g}\|d-d'\|_Z.
\]
The two coordinates in \eqref{eq:absorb-inverse} therefore have Lipschitz constants at most
\(1+L_Kq/c_g\) and \(q/c_g\), respectively. Their maximum is the
second estimate in \eqref{eq:absorb-bounds}, and its reciprocal is a lower Lipschitz
constant for \(F\).
\end{proof}

The remaining sections supply this criterion with \(E=c_0\). In
particular, the inverse calculation uses only \(QK=g-P\). It never
asks for an element \(b\) with prescribed \(Qb\) until the formula
\(b=d-K(g^{-1}(Qd))\) has already produced one.
